% frontmatter/czytam/como-usar.tex -- las instrucciones para el adulto de
% «Czytam z lupą», en polaco: las de frontmatter/lupa/como-usar.tex, con
% los nexos del polaco (że, bo, chociaż, kiedy, jeśli, żeby), las
% palabras del texto medidas en polaco (content/czytam/progresion.json)
% y de dónde viene la historia (el cuaderno en español, semana a semana).
% Lo que se dice del niño o de la niña no tiene género: «dziecko».
\section*{Jak korzystać z tego zeszytu}

To trzeci zeszyt z serii \emph{Uczę się czytać} po polsku, po
\emph{Pierwszych słowach} i \emph{Zdaniach}. Jest dla dziecka, które
czyta już płynnie całe zdania, także złożone -- na przykład dlatego, że
skończyło \emph{Zdania} -- i teraz robi następny krok: rozumie długie
zdania, z kilkoma częściami połączonymi słowami \textit{że},
\textit{bo}, \textit{chociaż}, \textit{kiedy}, \textit{jeśli} albo
\textit{żeby}, i zwraca uwagę na to, co tekst mówi naprawdę, a nie na
to, co wydaje się, że mówi. Czyta z lupą.

Tak jak w poprzednich zeszytach, na każdy dzień od poniedziałku do
piątku jest jedna strona, przez cały rok, także w czasie ferii i
wakacji, i każda strona wygląda zawsze tak samo:

\begin{itemize}[leftmargin=6mm, itemsep=1mm]
\item \textbf{Data}, \textbf{Czytanie} i \textbf{Uwagi}: dla dorosłego.
\item \textbf{Tekst dnia}, w zielonej ramce, w akapitach: około 45 słów
  na początku roku i około 110 na końcu. Tytuł ramki mówi, jak go czytać
  (na głos, dwa razy, głosem każdej postaci, najpierw po cichu...).
\item \textbf{Zadanie}, do którego trzeba wrócić do tekstu. Złota
  zasada: \textbf{odpowiedź zawsze jest w tekście} -- w tekście z tego
  dnia albo, jeśli tak mówi polecenie, w tekstach z poprzednich dni tego
  samego tygodnia. Jeśli dziecko jej nie znajduje, nie mów mu jej:
  zapytaj „gdzie to jest napisane?” i pozwól mu przeczytać jeszcze raz.
\end{itemize}

\textbf{Zadania.} \textbf{\lblDibujaConLupa}: narysować dokładnie to, co
opisuje tekst; nie chodzi o to, żeby rysunek był ładny, tylko żeby był
dokładny (ile, jakiego koloru, gdzie), a na koniec sprawdza się go z
tekstem w ręku, odpowiadając na pytania pod spodem. \textbf{\lblMapa}:
iść za wskazówkami po kratkach i zaznaczyć miejsce albo drogę.
\textbf{\lblLogica}: zagadka logiczna; w każdej kratce stawia się ✓ albo
✗, aż zostanie tylko jedna możliwość. \textbf{\lblErrores}: streszczenie
tekstu z błędami, które trzeba znaleźć i poprawić.
\textbf{\lblVerdaderoFalso}: a jeśli fałsz, to jak jest naprawdę.
\textbf{\lblCodigo}: odczytać wiadomość szyfrem, który wyjaśnia tekst.
\textbf{\lblResponde}, \textbf{\lblOrdena}, \textbf{\lblRelaciona},
\textbf{\lblFicha}, \textbf{\lblCompara}, \textbf{\lblAdivina},
\textbf{\lblVinetas} i \textbf{\lblCrea}: pytania „dlaczego?” i „skąd to
wiesz?”, zdarzenia do ułożenia po kolei, informacje do wypisania z
tekstu popularnonaukowego, dwie rzeczy do porównania, zagadki, historia
w trzech rysunkach i trochę własnego pisania.

\textbf{Co tydzień zagadka.} Strony opowiadają o kolejnym roku z życia
\textbf{Lucíi}, \textbf{Daniego}, \textbf{Toby'ego} i ich rodziny, i o
\textbf{Klubie Lupy}, klubie detektywów, który Lucía zakłada ze swoimi
przyjaciółmi, Sofíą i Hugiem. Prawie co tydzień jest tajemnica:
wskazówki pojawiają się od poniedziałku do czwartku, a w piątek trzeba
ją rozwiązać -- \textbf{\lblCaso}: zakreślić odpowiedź i wypisać
wskazówki, które jej dowodzą. Rozwiązanie potwierdza się w następny
poniedziałek, a jest też w \textbf{odpowiedziach} na końcu zeszytu.
Lepiej tam nie zaglądać, zanim dziecko samo spróbuje: są po to, żeby
sprawdzić, albo żeby dać małą podpowiedź, kiedy utknie. W każdej części
roku jest też jedna wielka tajemnica, która wyjaśnia się w jej ostatnim
tygodniu. To ta sama historia, tydzień po tygodniu, co w hiszpańskim
zeszycie \emph{Leo con lupa}; rodzina mieszka w Hiszpanii, więc są w
niej też hiszpańskie zwyczaje: Trzech Króli, dwanaście winogron w
sylwestra, karnawał.

\textbf{Tempo.} Dziesięć albo piętnaście minut dziennie. Jeśli przez
cały tydzień dziecko czyta i rozwiązuje strony bez pomocy, może robić
dwie dziennie; jeśli tekst wydaje mu się za długi, przeczytajcie go
najpierw razem na głos, po jednym akapicie, a potem niech przeczyta go
jeszcze raz samo. Najważniejsze, żeby to wciąż była zabawa: detektywi
się nie spieszą, detektywi są uważni.

Na koniec każdej części jest medal, a na końcu zeszytu -- dyplom.
\newpage
